\documentclass[11pt]{article}
\usepackage[margin=1in]{geometry}
\usepackage{xcolor}
\usepackage[breakable,listings]{tcolorbox}
\tcbuselibrary{listings}
\usepackage{hyperref}
\usepackage{enumitem}
\usepackage{listings}
\usepackage{amsmath}
\usepackage{booktabs}
\usepackage{array}

% Define colors
\definecolor{osaorange}{RGB}{220,115,38}
\definecolor{aiblue}{RGB}{30,144,255}

% Configure hyperref
\hypersetup{
    colorlinks=true,
    linkcolor=blue,
    urlcolor=blue
}

% Code environment
\newtcblisting{rcode}{
  colback=white,
  colframe=gray!50,
  boxrule=0.5pt,
  arc=0pt,
  left=3pt,
  right=3pt,
  top=3pt,
  bottom=3pt,
  width=\textwidth,
  breakable,
  listing only,
  listing options={
    basicstyle=\small\ttfamily,
    breaklines=true,
    breakatwhitespace=true,
    columns=flexible,
    keepspaces=true,
    showstringspaces=false
  }
}

\title{}
\author{}
\date{}

\begin{document}

% Header box
\begin{tcolorbox}[colback=osaorange,colframe=black,boxrule=2pt,arc=0pt,left=3pt,right=3pt,top=6pt,bottom=6pt,boxsep=2pt]
\centering
\textcolor{white}{\textbf{\Large H524 Introduction to Biostatistics}}\\
\textcolor{white}{\textbf{\Large Week 6: Power, Two-Sample Tests, and ANOVA}}\\
\textcolor{white}{\textbf{\Large Homework Assignment}}\\
\textcolor{white}{\textbf{\Large Fall 2025}}
\end{tcolorbox}

\vspace{8pt}

\noindent\textbf{Due Date:} End of Week 6\\
\textbf{Total Points:} 10\\
\textbf{Format:} Submit through Canvas\\
\textbf{Time Limit:} 60 minutes

\section*{Instructions}

\begin{itemize}
\item Show all work for full credit
\item Round numerical answers to 2 decimal places unless otherwise specified
\item Include R code where specified
\item You may use R to perform calculations
\end{itemize}

\section{Numerical Questions (10 points total)}

\subsection*{Question 1: Sample Size for Power (5 points)}

A researcher is planning a study to test whether a new medication reduces systolic blood pressure. Based on pilot data, the population standard deviation is estimated to be $\sigma = 20$ mmHg. The researcher wants to detect a mean reduction of 10 mmHg from the population mean of 145 mmHg.

\vspace{0.3cm}

Using a significance level of $\alpha = 0.05$ (two-sided) and desired power of 0.80, calculate the required sample size using R's \texttt{power.t.test()} function for a one-sample t-test.

\vspace{0.3cm}

\textbf{Show your R code:}

\vspace{3cm}

\textbf{What is the required sample size (round UP to the nearest whole number)?}

\vspace{0.3cm}

\textbf{Answer:} \underline{\hspace{3cm}} subjects

\vspace{1cm}

\subsection*{Question 2: Two-Sample t-Test and ANOVA (5 points)}

A researcher compares the effectiveness of three different exercise programs on weight loss (in kg) over 8 weeks. The data are:

\vspace{0.3cm}

\begin{center}
\begin{tabular}{lll}
\toprule
Program A & Program B & Program C \\
\midrule
4.2 & 6.1 & 7.8 \\
3.8 & 5.8 & 8.2 \\
4.5 & 6.4 & 7.5 \\
4.1 & 5.9 & 8.0 \\
4.3 & 6.2 & 7.9 \\
\bottomrule
\end{tabular}
\end{center}

\vspace{0.3cm}

\textbf{Part A:} Perform a one-way ANOVA to test if there is a significant difference in mean weight loss across the three programs. Use $\alpha = 0.05$.

\vspace{0.3cm}

\textbf{Show your R code to create the data and perform ANOVA:}

\vspace{5cm}

\textbf{What is the F-statistic (round to 2 decimal places)?}

\vspace{0.3cm}

\textbf{Answer:} F = \underline{\hspace{3cm}}

\vspace{0.5cm}

\textbf{Part B:} Compare ONLY Program A vs. Program C using a two-sample t-test (assume equal variances, two-sided test).

\vspace{0.3cm}

\textbf{Show your R code:}

\vspace{4cm}

\textbf{What is the p-value (round to 4 decimal places)?}

\vspace{0.3cm}

\textbf{Answer:} p-value = \underline{\hspace{3cm}}

\vspace{1cm}

\end{document}
